\documentclass[11pt]{article}
\usepackage{srmstyle}
\setsrmdoctitle{Guide technique}

\begin{document}

\srmtitlepage{Guide technique}{Manuel du developpeur SRM Fleet}{Comprendre, faire evoluer et maintenir l'application}

\tableofcontents
\vspace{0.6cm}
\srmrule
\vspace{0.4cm}

% =====================================================================
\section{Introduction}
\textbf{SRM Fleet} est le backoffice web de gestion du parc automobile de SRM Souss-Massa. Il couvre
le referentiel des vehicules et des conducteurs, le cycle de vie des affectations, les cartes
carburant, la conformite reglementaire (assurance, visite technique, sinistres), un moteur d'alertes,
un tableau de bord et des exports CSV et PDF.

Ce guide s'adresse a toute personne qui reprendra le code. Il explique l'architecture, le
fonctionnement d'un module de bout en bout, les briques communes a tous les modules, puis la facon
d'ajouter un module en respectant les conventions du projet. Il reste volontairement synthetique :
un seul module est detaille, les autres suivent le meme schema.

\begin{center}
\renewcommand{\arraystretch}{1.3}
\begin{tabularx}{\linewidth}{>{\bfseries}l Y}
\toprule
\rowcolor{srmsoft} Couche & Technologie \\
\midrule
Frontend & Angular 21 (composants standalone, signaux), design system maison, servi par Nginx. \\
Backend  & Spring Boot 3.5 sur Java 21, Spring Security (JWT), Spring Data JPA. \\
Base de donnees & PostgreSQL 17, migrations Flyway. \\
Generation PDF & OpenPDF. \\
Orchestration & Docker Compose. \\
\bottomrule
\end{tabularx}
\end{center}

% =====================================================================
\section{Vue d'ensemble de l'architecture}
L'application suit une architecture classique en trois tiers, sans etat cote serveur.

\begin{srmnote}
\textbf{Navigateur} (Angular) \;$\rightarrow$\; \textbf{Nginx} sert l'application et relaie \texttt{/api/}
\;$\rightarrow$\; \textbf{API REST Spring Boot} \;$\rightarrow$\; \textbf{PostgreSQL}.\\
L'authentification repose sur un jeton \textbf{JWT} envoye dans l'en-tete \texttt{Authorization}. Le
backend est \textbf{stateless} : aucune session serveur, chaque requete porte son jeton.
\end{srmnote}

\subsection{Deux roles, un cloisonnement}
\begin{itemize}
  \item \textbf{Administrateur central} : visibilite nationale, gere le referentiel et valide les demandes de carte.
  \item \textbf{Responsable provincial} : acces limite a sa province, sans visibilite sur les autres.
\end{itemize}
Ce cloisonnement (regle \textbf{RG14}) est applique dans la couche service du backend. Il est present
dans presque tous les modules : c'est un reflexe a garder en tete pour toute nouvelle fonctionnalite.

% =====================================================================
\section{Le backend}

\subsection{Une organisation par module}
Le code est range \textbf{par domaine metier}, pas par type technique. Chaque module contient sa propre
pile de classes, ce qui rend chaque fonctionnalite autonome et facile a retrouver. Exemple pour le
module Vehicule :
\begin{lstlisting}
ma/srm/fleet/referentiel/vehicule/
  VehiculeController.java       # exposition REST (couche mince)
  VehiculeService.java          # regles metier + transactions
  VehiculeRepository.java       # acces base (Spring Data JPA)
  VehiculeSpecifications.java   # filtres de recherche dynamiques
  Vehicule.java                 # entite JPA (table)
  StatutVehicule.java           # enumerations du domaine
  dto/VehiculeRequest.java      # donnees entrantes (validees)
  dto/VehiculeResponse.java     # donnees sortantes (jamais l'entite)
\end{lstlisting}

\subsection{Le parcours d'une requete}
Toutes les requetes suivent le meme trajet en couches. On le detaille ici sur la creation d'un
vehicule ; \textbf{tous les modules fonctionnent ainsi}.

\subsubsection{1. Le controleur : mince, il ne fait que router}
Il declare la route, controle le role si besoin, et delegue au service. Aucune logique metier ici.
\begin{lstlisting}[language=Java]
@PostMapping
@ResponseStatus(HttpStatus.CREATED)
@PreAuthorize("hasRole('ADMINISTRATEUR_CENTRAL')")
public VehiculeResponse creer(@Valid @RequestBody VehiculeRequest req) {
    return service.creer(req);
}
\end{lstlisting}
Points a retenir : \texttt{@Valid} declenche la validation des champs, \texttt{@PreAuthorize} restreint
l'acces au role, \texttt{@AuthenticationPrincipal UserPrincipal user} (sur les lectures) fournit
l'utilisateur connecte pour le cloisonnement.

\subsubsection{2. Le service : les regles metier et la transaction}
Annote \texttt{@Transactional}, il porte toute la logique : unicite, regles de gestion, resolution des
references, puis conversion en DTO.
\begin{lstlisting}[language=Java]
public VehiculeResponse creer(VehiculeRequest req) {
    if (vehiculeRepository.existsByImmatriculation(req.immatriculation())) {
        throw new BusinessException("Cette immatriculation existe deja", HttpStatus.CONFLICT);
    }
    Vehicule v = new Vehicule();
    appliquer(v, req);                       // RG2 : location => fournisseur obligatoire
    return toResponse(vehiculeRepository.save(v));
}
\end{lstlisting}

\subsubsection{3. Repository, entite et DTO}
Le \textbf{repository} etend \texttt{JpaRepository} : les methodes de base (save, findById, findAll) sont
fournies ; on ajoute des methodes derivees du nom (\texttt{existsByImmatriculation}). L'\textbf{entite}
represente la table. Le \textbf{DTO} (un \texttt{record} Java) transporte les donnees : l'entite JPA ne
quitte jamais le service, le controleur ne parle qu'en DTO.

\begin{srmnote}
Regle d'or du projet : \textbf{injection par constructeur}, \textbf{aucune logique metier dans les
controleurs}, \textbf{les entites ne sortent pas de la couche service}. Le mapping entite vers DTO est
ecrit a la main (methode \texttt{toResponse} ou classe \texttt{...Mapper}), sans librairie.
\end{srmnote}

\subsection{Les memes couches partout}
Chaque module reprend ce schema. La liste ci-dessous situe les principaux modules ; leur structure
interne est identique a celle du module Vehicule.
\begin{center}
\renewcommand{\arraystretch}{1.25}
\begin{tabularx}{\linewidth}{>{\bfseries}l Y}
\toprule
\rowcolor{srmsoft} Module & Role \\
\midrule
Referentiel & Vehicules, conducteurs et permis, fournisseurs, provinces, departements. \\
Affectation & Lien vehicule vers conducteur, historique, cloture. \\
Carte carburant & Demande, validation par le central, emission de la carte. \\
Conformite & Assurances, visites techniques, sinistres et pieces jointes. \\
Alerte & Moteur planifie et configuration des seuils. \\
Dashboard & Agregation des statistiques de tous les modules. \\
Export & Generation CSV et PDF, reutilisant les filtres des listes. \\
Admin et Auth & Comptes utilisateurs, connexion et jetons JWT. \\
\bottomrule
\end{tabularx}
\end{center}

\subsection{La securite}
La configuration est centralisee dans \texttt{SecurityConfig}. Choix structurants : \textbf{CSRF
desactive} (inutile pour une API a jeton), \textbf{sessions STATELESS}, et une \textbf{liste blanche}
d'endpoints publics (login, healthcheck, documentation Swagger). Tout le reste exige une
authentification (principe du moindre privilege).
\begin{itemize}
  \item \textbf{Connexion} : \texttt{/api/v1/auth/login} verifie les identifiants via
        l'\texttt{AuthenticationManager}, puis \texttt{JwtService} genere le jeton.
  \item \textbf{Chaque requete suivante} : le \texttt{JwtAuthenticationFilter} lit le jeton, le valide et
        peuple le contexte de securite avant les filtres standards de Spring.
  \item \textbf{Roles} : \texttt{@PreAuthorize("hasRole('ADMINISTRATEUR\_CENTRAL')")} protege les actions
        reservees au central (creation de referentiel, validation des demandes).
  \item \textbf{Mots de passe} : haches avec \textbf{BCrypt} (sel integre).
\end{itemize}

\subsection{Le cloisonnement par province (RG14)}
Dans les services, un Responsable provincial ne voit que sa province. En \textbf{lecture de liste}, on
ajoute une specification qui filtre sur \texttt{provinceId}. En \textbf{acces a un element precis}, on
verifie la portee ; si l'element appartient a une autre province, on renvoie \textbf{404} (et non 403)
pour ne pas reveler son existence.
\begin{lstlisting}[language=Java]
if (user.getRole() == RoleUtilisateur.RESPONSABLE_PROVINCIAL
        && !v.getProvince().getId().equals(user.getProvinceId())) {
    throw new BusinessException("Vehicule introuvable : " + v.getId(), HttpStatus.NOT_FOUND);
}
\end{lstlisting}

\subsection{Les filtres de recherche (Specifications)}
Les listes acceptent des filtres optionnels (immatriculation, type, statut, province). On les compose
avec l'API \texttt{Specification} de Spring Data : chaque critere est une brique, ignoree si le
parametre est absent. Le cloisonnement province se branche comme une brique de plus.

\subsection{Validation et gestion des erreurs}
Deux niveaux, une reponse homogene :
\begin{itemize}
  \item \textbf{Validation des champs} : annotations Bean Validation sur les DTO (\texttt{@NotBlank},
        \texttt{@NotNull}...), declenchees par \texttt{@Valid}.
  \item \textbf{Regles metier} : on leve une \texttt{BusinessException} portant un message et un statut HTTP.
\end{itemize}
Le \texttt{GlobalExceptionHandler} (annote \texttt{@RestControllerAdvice}) transforme toute exception en
reponse \textbf{ProblemDetail} (norme RFC 7807 : un corps JSON standard avec titre, statut, detail). Les
erreurs inattendues renvoient un 500 generique, sans exposer de detail technique.

\subsection{Les cycles de vie}
Plusieurs modules portent une petite machine a etats, geree dans le service (transitions verifiees,
statut stocke en base) :
\begin{itemize}
  \item \textbf{Demande de carte} : \texttt{EN\_ATTENTE} puis \texttt{VALIDEE} ou \texttt{REFUSEE}. La
        validation (par le central) fixe le plafond et le fournisseur, et emet la carte.
  \item \textbf{Affectation} : \texttt{ACTIVE} puis \texttt{CLOTUREE}. La cloture conserve l'historique.
  \item \textbf{Sinistre} : \texttt{OUVERT}, \texttt{EN\_COURS}, \texttt{CLOTURE}, avec pieces jointes.
\end{itemize}

\subsection{Le moteur d'alertes}
Un service planifie (Spring Scheduler, une fois par jour) parcourt les configurations de seuils. Pour
chaque type surveille (assurance, visite technique), il retient les echeances qui tombent dans la
fenetre du seuil et cree une alerte \textbf{si elle n'existe pas deja} (deduplication par type et
reference). Les notifications sont \textbf{in-app} uniquement (canal tableau de bord) ; le tableau de
bord des alertes les classe par urgence.

\subsection{Base de donnees et migrations}
Le schema est gere par \textbf{Flyway} : des fichiers \texttt{V1\_\_...sql} a \texttt{V8\_\_...sql}
appliques automatiquement au demarrage, dans l'ordre. On ne modifie jamais une migration deja livree :
on en ajoute une nouvelle.
\begin{srmwarn}
Deux regles critiques reposent sur des \textbf{index uniques partiels PostgreSQL}, impossibles a
reproduire avec une base en memoire : \textbf{RG3} (une seule affectation active par vehicule) et
\textbf{RG4} (une seule carte active par affectation). C'est la base qui garantit ces regles, en plus du
controle applicatif. Elles sont verifiees par des tests d'integration Testcontainers.
\end{srmwarn}
Le \textbf{seed} initial insere le referentiel de base et le compte \texttt{admin} (dont le mot de passe
est pose au demarrage par \texttt{AdminPasswordInitializer}, a partir de la configuration : aucun mot de
passe par defaut, voir le guide de deploiement).

\subsection{Les exports CSV et PDF}
Chaque liste propose un bouton d'export. Les exports \textbf{reutilisent les services de liste} : memes
filtres et meme cloisonnement province que l'ecran. Le \textbf{CSV} est genere sans dependance (UTF-8
avec BOM pour Excel, separateur point-virgule). Le \textbf{PDF} (OpenPDF) est mis en page selon la charte
graphique : logo, filet tricolore officiel, tableau a en-tete bleu, pied de page.

% =====================================================================
\section{Le frontend}

\subsection{Organisation}
L'application Angular est decoupee en trois espaces :
\begin{itemize}
  \item \texttt{core/} : services transverses (authentification, theme, composants d'interface communs).
  \item \texttt{features/} : un dossier par module (vehicules, conducteurs, affectations...), avec ses
        composants, son service HTTP, ses modeles et ses routes.
  \item \texttt{layout/} : la coquille de l'application (menu lateral, barre superieure).
\end{itemize}
Tous les composants sont \textbf{standalone} (pas de NgModule), l'etat local s'exprime avec des
\textbf{signaux}, et le typage est strict (pas de \texttt{any}).

\subsection{Le parcours d'un module (cote client)}
Meme regularite qu'au backend. Pour un module donne :
\begin{itemize}
  \item un \textbf{service} expose les appels HTTP vers \texttt{/api/v1/...} et renvoie des \texttt{Observable} ;
  \item trois \textbf{composants} type : la \textbf{liste} (recherche et filtres), le \textbf{formulaire}
        (creation et modification), le \textbf{detail} (fiche en lecture) ;
  \item les \textbf{modeles} (interfaces TypeScript) decrivent les donnees echangees.
\end{itemize}
\begin{lstlisting}
rechercherVehicules(f: RechercheVehicule): Observable<Vehicule[]> {
  let params = new HttpParams();
  if (f.immatriculation) params = params.set('immatriculation', f.immatriculation);
  if (f.statut) params = params.set('statut', f.statut);
  return this.http.get<Vehicule[]>('/api/v1/vehicules', { params });
}
\end{lstlisting}

\subsection{L'authentification cote client}
Le \texttt{AuthService} est la source de verite de la session. Il conserve le jeton et le profil dans le
\texttt{localStorage} et expose des signaux (\texttt{isAuthenticated}, \texttt{isAdmin}). Deux mecanismes
completent le tableau :
\begin{itemize}
  \item \textbf{Intercepteurs HTTP} : \texttt{authInterceptor} ajoute l'en-tete \texttt{Authorization} sur
        les appels \texttt{/api/} ; \texttt{errorInterceptor} traite les erreurs (deconnexion sur 401).
  \item \textbf{Gardes de routes} : \texttt{authGuard} (routes authentifiees), \texttt{guestGuard} (page de
        connexion), \texttt{adminGuard} (routes reservees au central).
\end{itemize}
La session est restauree de facon synchrone au demarrage, avant toute decision de routage. La validite
du jeton (champ \texttt{exp}) est verifiee cote client par confort ; la validation qui fait foi reste
celle du backend.

\subsection{Le design system maison}
Aucune librairie d'interface externe. Tout part des \textbf{jetons} (couleurs, espacements, rayons,
typographie) et des \textbf{primitives} CSS definies dans \texttt{frontend/src/styles.scss} :
\texttt{.card}, \texttt{.btn}, \texttt{.field}, \texttt{.input}, \texttt{.table}, \texttt{.badge}...
Chaque nouvel ecran se compose a partir de ces primitives, ce qui garantit une coherence visuelle totale
et un theme clair et sombre gere en un seul endroit. Les graphes du tableau de bord sont dessines en SVG.

% =====================================================================
\section{Les regles de gestion cles}
Ces regles structurent le metier. Elles sont appliquees dans les services et, pour les plus critiques,
garanties par la base.
\begin{center}
\renewcommand{\arraystretch}{1.3}
\begin{tabularx}{\linewidth}{>{\bfseries}l Y}
\toprule
\rowcolor{srmsoft} Regle & Enonce \\
\midrule
RG2  & Un vehicule en location doit referencer un fournisseur. \\
RG3  & Un vehicule ne peut avoir qu'une seule affectation active a la fois (index unique partiel). \\
RG4  & Une affectation ne peut avoir qu'une seule carte carburant active (index unique partiel). \\
RG11 & Le conducteur accompagnateur d'une visite est deduit de l'affectation active a la date. \\
RG14 & Un Responsable provincial n'accede qu'aux donnees de sa province. \\
\bottomrule
\end{tabularx}
\end{center}

% =====================================================================
\section{Ajouter un nouveau module}
La regularite du projet rend l'ajout d'un module mecanique. Voici la recette, en reprenant le module
Vehicule comme modele.

\subsection{Cote backend}
\begin{enumerate}
  \item Creer le package \texttt{ma/srm/fleet/<domaine>/<module>/}.
  \item Ecrire l'\textbf{entite} JPA et l'ajouter au schema via une \textbf{nouvelle} migration Flyway
        (\texttt{V9\_\_...sql}), sans toucher aux migrations existantes.
  \item Definir les \textbf{DTO} (record \texttt{...Request} valide, \texttt{...Response}).
  \item Creer le \textbf{repository} (\texttt{extends JpaRepository}).
  \item Ecrire le \textbf{service} (\texttt{@Transactional}) : regles metier, cloisonnement province,
        mapping vers DTO ecrit a la main.
  \item Ajouter les \textbf{Specifications} si le module a une recherche filtrable.
  \item Exposer le \textbf{controleur} REST (\texttt{/api/v1/...}), mince, avec \texttt{@PreAuthorize} la
        ou c'est necessaire.
  \item Ajouter des \textbf{tests} (unitaires, et d'integration Testcontainers pour toute contrainte de base).
\end{enumerate}

\subsection{Cote frontend}
\begin{enumerate}
  \item Creer le dossier \texttt{features/<module>/}.
  \item Definir les \textbf{modeles} TypeScript et le \textbf{service} HTTP vers \texttt{/api/v1/...}.
  \item Ecrire les composants \textbf{liste}, \textbf{formulaire}, \textbf{detail}, a partir des
        \textbf{primitives} du design system (aucun style redefini inutilement).
  \item Declarer les \textbf{routes} du module et les rattacher au menu, avec les gardes adaptees.
\end{enumerate}

\begin{srmok}
Regle simple : un nouveau module doit ressembler aux modules existants. En cas de doute, copier la
structure du module Vehicule (backend) et du module Vehicules (frontend), puis adapter.
\end{srmok}

% =====================================================================
\section{Conventions Git et flux de travail}

\subsection{Branches}
\begin{itemize}
  \item \texttt{main} : code livrable, protegee.
  \item \texttt{develop} : integration continue, protegee.
  \item \texttt{feature/<module>-<description>} : une fonctionnalite.
  \item \texttt{fix/<module>-<description>} : une correction. \texttt{hotfix/<description>} : correction urgente sur \texttt{main}. \texttt{chore/<description>} : tache technique.
\end{itemize}

\subsection{Commits (Conventional Commits)}
Format : \texttt{type(scope): description a l'imperatif}. Le \texttt{scope} est le module concerne.
\begin{center}
\renewcommand{\arraystretch}{1.25}
\begin{tabularx}{\linewidth}{>{\ttfamily}l Y}
\toprule
\rowcolor{srmsoft} \textmd{\bfseries Type} & Usage \\
\midrule
feat     & Nouvelle fonctionnalite. \\
fix      & Correction de bug. \\
refactor & Reorganisation sans changement de comportement. \\
test     & Ajout ou modification de tests. \\
docs     & Documentation. \\
chore, build, ci, perf & Taches techniques, outillage, integration continue, performance. \\
\bottomrule
\end{tabularx}
\end{center}
Exemple : \texttt{feat(carte): validation d'une demande et emission de la carte}.

\subsection{Pull requests}
\begin{enumerate}
  \item Rebaser sur \texttt{develop} a jour avant d'ouvrir la PR.
  \item Une PR egale une tache, petite et ciblee ; remplir le modele (objectif, perimetre, test).
  \item CI verte obligatoire ; un relecteur obligatoire (l'autre developpeur), pas d'auto-fusion.
  \item Fusion en \textbf{squash} vers \texttt{develop}, puis suppression de la branche.
\end{enumerate}

\subsection{Regles de code}
\begin{itemize}
  \item Backend : injection par constructeur, \texttt{@Transactional} en couche service, aucune logique dans les controleurs.
  \item Frontend : composants standalone, typage strict, pas de \texttt{any}.
  \item Formatage verifie en CI : Spotless et Google Java Format cote backend, ESLint et Prettier cote frontend.
\end{itemize}

% =====================================================================
\section{Decisions d'architecture (resume)}
Les choix structurants sont consignes sous forme de decisions courtes. Resume :
\begin{center}
\renewcommand{\arraystretch}{1.3}
\begin{tabularx}{\linewidth}{>{\bfseries}l Y}
\toprule
\rowcolor{srmsoft} Decision & Raison \\
\midrule
Design system maison & Coherence visuelle totale, pas de librairie d'interface externe. \\
Mappers manuels & Mapping explicite et lisible, sans generation de code (pas de MapStruct). \\
Alertes in-app uniquement & Pas de service courriel valide en production ; canal courriel reporte. \\
Exports CSV et PDF & OpenPDF aux couleurs de la charte, un seul generateur pour tous les modules. \\
Tests Testcontainers & Valider sur un vrai PostgreSQL les regles portees par la base (RG3, RG4). \\
Init admin securisee & Aucun mot de passe par defaut : le compte admin vient de la configuration. \\
\bottomrule
\end{tabularx}
\end{center}

% =====================================================================
\section{Tests et integration continue}
\begin{itemize}
  \item \textbf{Backend} : tests unitaires et tests d'integration Testcontainers (PostgreSQL reel via
        Docker), qui valident les migrations, le seed et les contraintes critiques RG3 et RG4.
  \item \textbf{Frontend} : \texttt{lint} puis \texttt{build} de production.
  \item \textbf{CI} (GitHub Actions) : sur chaque push et pull request vers \texttt{main} et
        \texttt{develop}, deux jobs (backend et frontend) doivent passer au vert avant toute fusion.
\end{itemize}

\vspace{0.4cm}
\srmrule
\begin{center}\small\textcolor{srmslate}{Fin du document. SRM Fleet, Oussama BENAZZOUZ et TAKI Fadlouallah.}\end{center}

\end{document}
